\section{Leech格提升中的赤道扩充}
\label{sec:equatorial}

码构造的局部替换需要在删除与插入之间取得净收益。本章有另一种机会：
一个大型提升构型留出了尚未使用的赤道。研究的目标是找到能够直接加入的方向，
使它们彼此相容，又同时满足整个提升层的约束。144条这样的直线贡献288个点，
将下界从591612提高到591900。

\subsection{提升构型与赤道约束}

记 \(\Sigma\subset\mathbb Z^{24}\) 为平方范数32的 Leech
极小壳，共196560个向量。它由所给扩展 Golay
码生成：1104个向量在两个坐标取
\(\pm4\)；97152个向量在一个八重量码字的支撑上取
\(\pm2\)，负号数为偶数；98304个向量先按码字取符号
\((-1)^{c_i}\)，再任选一个坐标反号并乘以3。Golay
生成矩阵秩12，重量分布为
\(1+759z^8+2576z^{12}+759z^{16}+z^{24}\)。该标准模型中不同壳向量的内积至多16
\cite{conway1999sphere}。

38维写成
\(\mathbb R^{24}\oplus\mathbb R^{14}\)。公开的591,612点构型已经把全部
Leech 极小方向用于混合层
\cite{kissingnumbers}，但这并不占满几何上的赤道空间。新的问题是寻找另一组方向，使它们与整个混合层相容。

先说明所用母构型。把 \(\Sigma\)
的98,280条反极直线分为644个类，每类代表满足 \(|u\cdot u'|\le8\)。在
\(\mathbb R^{14}\) 中取1932个平方范数为8的方向
\(d\)，不同方向内积至多4，并把它们分成644个零和三角形。每个三角形内两两内积为
\(-4\)。将一类 Leech
直线配给一个尾三角形，每条直线取两种符号、三种尾标签，得到六个点

\[
\left(\frac{\varepsilon u}{\sqrt{48}},\frac d{\sqrt{24}}\right),
\qquad \varepsilon\in\{1,-1\}.
\]

这些点的平方范数是 \(32/48+8/24=1\)。同一头方向、不同尾标签的内积为
\(32/48-4/24=1/2\)；不同头方向、相同尾标签的内积至多
\(8/48+8/24=1/2\)。不同类的点对则由 \(16/48+4/24=1/2\)
控制；其余情形更小。因此混合层包含 \(6\cdot98280=589680\) 个点。

同一直线的相反头方向也包含在计数中：头部内积为 $-32/48$，即使尾标签相同，
总内积也只有 $-1/3$，所以两种符号可以同时保留。

纯尾层使用1932个点 \((0,Rd/\sqrt8)\)。数据给出的旋转为 \(R=Q/D\)，其中
\(Q\) 为整数矩阵、\(D>0\) 为整数，且

\[
QQ^{\mathsf T}=D^2I,
\qquad |d\cdot Qd'|\le\sqrt{48}\,D.
\]

第一式保持纯尾层的内积，第二式保证纯尾与混合层的内积不超过
\(1/2\)。右侧根式可通过整数平方根精确比较。由此得到基准点数
\(589680+1932=591612\)。

\subsection{第二壳层的相容直线}

采用同一个 Leech 格中平方范数48的壳层。
设 $\mathcal G$ 为所给 Golay 码，其整数坐标模型为
\begin{equation}
\mathcal L=\left\{x\in\mathbb Z^{24}:\begin{array}{l}
x_i\equiv m\pmod2\quad(m\in\{0,1\}),\\
((x_i-m)/2\bmod2)_{i=1}^{24}\in\mathcal G,\\
\sum_i x_i\equiv4m\pmod8
\end{array}\right\}.
\label{eq:leech-membership}
\end{equation}
由 $\mathcal G$ 的线性及全一码字，这个集合对加减封闭，因而是一个格。
其最小平方范数为32。偶坐标情形中，非零 Golay 剩余字至少有八个非零位，
对应坐标绝对值至少为二；若剩余字为零，则各坐标都是四的倍数，坐标和条件
排除了只有一个坐标绝对值为四的情况。奇坐标情形中，平方范数小于32的向量
只能各位取 $\pm1$；Golay 码重量为四的倍数，使坐标和模八为零，
与最后一个同余条件矛盾。前述极小向量达到平方范数32。

\begin{lemma}[第二壳层的相容性]
\label{lem:second-shell}
若 $v,u\in\mathcal L$ 的平方范数分别为48与32，则 $|v\cdot u|\le24$。
\end{lemma}
\begin{proof}
范数不同保证 $v\pm u\ne0$，故
$\|v\pm u\|^2=80\pm2v\cdot u\ge32$，即得所需内积界。
\end{proof}

母构型使用了全部极小直线，但赤道点可以来自不同长度的方向。
选择平方范数48而非32的整数向量后，归一化分母与混合层头部相同；
新的相容条件因此成为阈值为24的整数内积检查。
所给144个向量均满足式~\eqref{eq:leech-membership}的三个条件。
再检查它们彼此的内积，得到

\[
\|v_i\|^2=48,\qquad
|v_i\cdot v_j|\le24\ (i\ne j),\qquad
|v_i\cdot u|\le24\ (u\in\Sigma).
\]

\begin{proposition}[赤道补充]
\label{prop:equatorial}
向混合层与纯尾层加入288个点 $(\pm v_i/\sqrt{48},0)$，得到38维的
591900点接吻构型。
\end{proposition}
\begin{proof}
将288个赤道点 \((\pm v_i/\sqrt{48},0)\) 加入上述构型。它们彼此的内积至多
\(1/2\)，与混合层的内积为
\(\pm v_i\cdot u/48\le1/2\)，与纯尾点正交。基准构型没有赤道点，因而不需删除任何旧点。于是

\[
K(38)\ge591612+2\cdot144=591900.
\]
\end{proof}

引理~\ref{lem:second-shell}使整个旧壳层不再成为搜索约束。生成满足格成员同余
条件的范数48整数向量，每条反极直线只取一个规范代表；两条代表若满足
$|v\cdot v'|>24$，就在它们之间连边。剩下的问题是这个有限图中的独立集选择：
每条选中直线恰好贡献两个新赤道点。联合替换删除的是新增直线在旧选集中的
邻点并集，因而即使无法直接插入一条直线，也可能通过同时替换多条直线取得净增益。
每次接受的选集均须满足内部相容性。验证程序还直接重算新增向量与整个
$\Sigma$ 的内积，独立核对跨壳引理在所给数据上的结论；组合选择过程中不必
重复进行这一计算。

这里有一个直接的几何解释。用 $C=\Sigma/\sqrt{32}$ 表示单位母壳，
新赤道方向 $z=v/\sqrt{48}$ 必须满足
$|z\cdot c|\le\sqrt{3/8}$。这个阈值严格大于 $1/2$，因为混合层的头部
只有 $\sqrt{2/3}$ 的长度。提升为赤道留下的空间，因而大于母壳自身允许
直接插入单位点的空间；第二壳层正是利用这部分空间。这也解释了为何
使用全部极小方向并不意味着相应的38维构型已经无法增加点。

搜索先得到142条相容直线，随后扩展为144条。最终数据与母构型使用同一份
Golay
生成矩阵；规范代表取首个非零坐标为正，再按整数元组的字典序排列。固定这一约定，使644类划分、尾三角形和新直线的点积检查彼此一致。
